% ====================================================
% 6. System Design
% ====================================================

\subsection{System Architecture}

Dikaios Amanah will follow a layered web application architecture designed
to separate user interaction, application processing, persistent data
management, document storage, and trust and audit functions.

The major components of the proposed architecture are:

\begin{itemize}
    \item Presentation Layer
    \item Application and API Layer
    \item Database Layer
    \item Document Storage
    \item Authentication and Security Services
    \item Permissioned Trust and Audit Layer
\end{itemize}

The overall architecture is intended to provide separation of
responsibilities while allowing the individual components to communicate
through defined interfaces.

The high-level architecture can be represented as follows:

\begin{figure}[H]
    \centering
    \fbox{
        \parbox{0.85\textwidth}{
            \centering
            \textbf{System Architecture Diagram Placeholder}\\[0.3cm]
            React Frontend $\rightarrow$ Node/Express Backend
            $\rightarrow$ PostgreSQL Database\\
            Backend $\rightarrow$ Firebase Storage\\
            Backend $\rightarrow$ Trust and Audit Layer
        }
    }
    \caption{High-Level System Architecture}
    \label{fig:system_architecture}
\end{figure}

The final architecture diagram will be added to the SRS after the
individual system components and data relationships have been finalized.


\subsection{Presentation Layer}

The Presentation Layer provides the user-facing web interface of
Dikaios Amanah.

The planned frontend will use React and will provide role-specific
interfaces for:

\begin{itemize}
    \item Business Owners
    \item Investors
    \item Administrators
\end{itemize}

The Presentation Layer will be responsible for displaying information,
collecting user input, performing client-side validation where appropriate,
and communicating with backend services through defined APIs.

Business logic and security-sensitive authorization decisions shall not
depend solely on client-side controls.


\subsection{Application and API Layer}

The Application Layer will provide the primary business logic of the
system.

The planned backend will use Node.js with Express.js to provide RESTful
API services.

The Application Layer will be responsible for:

\begin{itemize}
    \item Processing client requests.
    \item Enforcing application business rules.
    \item Authenticating users.
    \item Authorizing role-specific operations.
    \item Managing business workflows.
    \item Processing investment-related operations.
    \item Coordinating database operations.
    \item Coordinating document storage operations.
    \item Recording critical audit events.
\end{itemize}

The backend will act as the central coordination layer between the
frontend and the system's persistent and supporting services.


\subsection{Database Layer}

PostgreSQL will be used as the primary relational database.

The database will maintain persistent application information including:

\begin{itemize}
    \item User accounts.
    \item Business information.
    \item Investor information.
    \item Verification records.
    \item Investment campaigns.
    \item Investment records.
    \item Profit information.
    \item Activity records.
    \item References to uploaded documents.
\end{itemize}

The database design will maintain appropriate relationships and integrity
constraints between related entities.

The detailed database structure and Entity Relationship Diagram will be
specified in the Data Requirements section.


\subsection{Document Storage Component}

Business verification and related processes may require supporting
documents.

The proposed system will use Firebase Storage as the designated document
storage mechanism.

The application database will store the relevant metadata and references
required to identify uploaded documents, while the actual document files
will be maintained by the designated storage service.

The system will generate cryptographic integrity information for applicable
documents so that authorized users can verify whether the document content
has been altered.


\subsection{Authentication and Authorization}

Authentication will establish the identity of a user before protected
system resources are accessed.

The proposed system will use token-based authentication, with JWT planned
for authenticated communication between the frontend and backend.

Authorization will be implemented through Role-Based Access Control
(RBAC).

The three primary roles are:

\begin{enumerate}
    \item Business Owner
    \item Investor
    \item Administrator
\end{enumerate}

Role verification shall occur at the backend level before protected
operations are executed.

This prevents users from gaining access to restricted operations simply
by manipulating the frontend interface.


\subsection{Permissioned Trust and Audit Layer}

Dikaios Amanah includes a Permissioned Trust and Audit Layer as an
additional mechanism for improving transparency and traceability.

The purpose of this layer is not to replace the primary PostgreSQL
database. Instead, it is intended to maintain tamper-evident records of
selected critical events.

Events that may require audit recording include:

\begin{itemize}
    \item Business verification.
    \item Contract approval.
    \item Investment confirmation.
    \item Profit declaration.
    \item Document integrity verification.
    \item Other critical administrative activities.
\end{itemize}

The trust layer is intended to provide an additional level of confidence
that important historical events have not been silently altered.

The MVP will implement the trust and audit concept at an appropriate level
for academic demonstration. Production-scale distributed ledger
deployment is outside the defined MVP scope.


\subsection{System Component Interaction}

The major system components will interact through the following general
workflow:

\begin{enumerate}
    \item A user interacts with the React-based frontend.
    \item The frontend sends a request to the backend API.
    \item The backend authenticates and authorizes the request.
    \item The backend applies the appropriate business rules.
    \item Required information is retrieved from or stored in PostgreSQL.
    \item If documents are involved, the backend communicates with the
    designated document storage service.
    \item If the operation is a critical auditable event, the appropriate
    audit information is recorded.
    \item The backend returns the result to the frontend.
    \item The frontend presents the result to the user.
\end{enumerate}


\subsection{Deployment Architecture}

The planned deployment architecture will separate the public frontend,
backend application services, and persistent data services.

The frontend may be deployed using a platform such as Vercel, while the
backend may be deployed on a suitable cloud server or VPS. PostgreSQL and
document storage will be provided through appropriate persistent cloud
services.

A conceptual deployment structure is:

\begin{figure}[H]
    \centering
    \fbox{
        \parbox{0.85\textwidth}{
            \centering
            \textbf{Deployment Diagram Placeholder}\\[0.3cm]
            User Browser $\rightarrow$ Frontend Hosting
            $\rightarrow$ Backend Server\\
            Backend Server $\rightarrow$ PostgreSQL\\
            Backend Server $\rightarrow$ Document Storage\\
            Backend Server $\rightarrow$ Trust and Audit Layer
        }
    }
    \caption{Conceptual Deployment Architecture}
    \label{fig:deployment_architecture}
\end{figure}

The final deployment diagram will be produced after the detailed
deployment configuration has been finalized.


\subsection{Design Principles}

The system design will follow the following principles:

\begin{itemize}
    \item Separation of frontend and backend responsibilities.
    \item Separation of application data and document storage.
    \item Backend-controlled authorization.
    \item Secure handling of authentication credentials.
    \item Clear separation between operational data and audit records.
    \item Modular design to support future expansion.
    \item Traceability between requirements, implementation, and testing.
\end{itemize}